\documentclass[12pt,letterpaper, onecolumn]{exam}
\usepackage{amsmath}
\usepackage{amssymb}
\usepackage{graphicx}
\usepackage{hyperref}
\thispagestyle{empty}

\begin{document}
\begingroup  
    \centering
    \LARGE Applications of Linear Algebra\\
    \LARGE Math 235H\\[0.5em]
    \large December 10, 2024\\[0.5em]
    \large Kavin Kannan\par
    \large Physics Major\par
\endgroup
\rule{\textwidth}{0.4pt}

\begingroup %Introduction
Physics is ultimately the building blocks of all other sciences. What mathematics is to 
physics: physics is to chemistry, biology, and engineering. This paper will focus on the 
former--more specifically, what linear algebra does for physics. Linear Algebra consists of a 
vast and powerful field of information that, when applied to physics, can revolutionize our 
understanding of it.\\
\indent Linear Algebra can help significantly simplify very complex systems of equations and, through 
vectors and matrices, can allow for the modeling of intricate physical systems to help notice 
or predict patterns in nature. The language of lienar algebra is intrinsic to physics as we 
label almost every value generated with as a "scalar" or a "vector". Matrices are critical to easily 
solving the countless linear systems generated by phyiscal phenomena. Eigenvalues 
give meaning to the Hamiltonian matrix. Higher Dimensional Matrices are our way to digest 
the incredibly complicated tensors that are applied to almost every field in physics. 
\endgroup
\rule{\textwidth}{0.4pt}

\begingroup %Linear Systems and Kirchoff's Laws
Linear Systems are critical to the daily functions of life, yet go largely unnoticed when
compared to more grandiose applications of linear algebra. The application of linear systems 
that I am intimately familiar with is Kirchoff's Laws. These laws, when applied to circuits 
consisting of resistors, capacitors, and inductors, can simplify the complex calculations for 
voltage, current, and resistance. These three quantities are some of the most important for 
day-to-day life as they define how the electricity, which defines our lives, actually works.\\
\indent Kirchoff's first law is that "the current flowing into a node (or a junction) must be equal 
to the current flowing out of it."\textsuperscript{1} Kirchoff's second law is that "in any complete loop within 
a circuit, the sum of all voltages across components which supply electrical energy (such as 
cells or generators) must equal the sum of all voltages across the other components in the 
same loop."\textsuperscript{1} These two laws in tandem create a system of linear equations involving three 
separate "loop currents" as unknowns.
\endgroup
\begin{figure}[h]
    \centering
    \begin{minipage}{0.45\textwidth}
        \centering
        \includegraphics[width=0.7\textwidth]{/Users/kavinkannan/Documents/Images/Kirchoff.png}
        \caption{Kirchoff's Law Circuit\textsuperscript{2}}
        \label{fig:image1}
    \end{minipage} \hfill
    \begin{minipage}{0.45\textwidth}
        \centering
        \includegraphics[width=0.9\textwidth]{/Users/kavinkannan/Documents/Images/Matrix.png}
        \caption{Kirchoff's Law Matrix\textsuperscript{2}}
        \label{fig:image2}
    \end{minipage}
    \end{figure}
\vspace{1em}

\newpage
\thispagestyle{empty}
\begingroup
Using the circuit above, we can create the matrix (also above) by considering each "loop" 
of the circuit to create its own Kirchoff's Law Equations. Using Ohm's Law (V=IR), we can 
easily obtain that:\\
15I\textsubscript{A}-5I\textsubscript{B}=15,\indent\indent\indent\indent\indent\indent\indent
-5I\textsubscript{A}+9I\textsubscript{B}-1\textsubscript{C}=0,\indent\indent\indent\indent\indent\indent\indent
-1I\textsubscript{B}+11\textsubscript{C}=0,\\
which translates easily into the matrix above. This matrix is easily solvable using elementary 
row operations (another linear algebra topic), and can thus solve for the currents in each 
loop. This information can then be used in understanding home circuitry, setting up power 
\endgroup
\rule{\textwidth}{0.4pt}

\begingroup %Eigenvectors and Eigenvalues
Eigenvectors and Eigenvalues can initially seem very redundant or insignificant, however, 
when observing almost every field, they always have some purpose. The Eigenvalue, 
in particular, when applied to physics, has a special purpose in regards to the energy spectrum 
of a system. The Hamiltonian operator is an operator "that represents the total energy of a 
quantum system, including both kinetic and potential energy."\textsuperscript{3} This operator acts on "the 
state space of a quantum system"\textsuperscript{3}, which is an abstract space representing not physical 
positions, but "quantum states" of the system.\\
\indent The Eigenvalues of the Hamiltonian "correspond to the possible energy levels of the 
system"\textsuperscript{3} and the Eigenvectors of those Eigenvalues "represent the corresponding stationary 
states"\textsuperscript{3}. If we were to apply the Hamiltonian to Shrödinger's Wave Equation, we can 
calculate the allowed energy levels and wave functions of a system:\\
$i\hbar \frac{d\psi}{dt} = H\psi$ where $H\psi = E\psi$ and $E$ is the total energy of the system.\\
The Hamiltonian Operator will eventually become $\frac{\hbar^{2}}{2m}\frac{\partial^{2}}{\partial t^{2}} + \mathsf{V}(t)$.\textsuperscript{4}\\
Ultimately, with a little more understanding of Quantumn Physics, the energy E of the 
system is equivalent to the "Eigenvalue of $\psi$ when the Hamiltonian acts on $\psi$."\textsuperscript{4}
\endgroup
\vspace{1em}
\begin{figure}[h]
    \centering
    \begin{minipage}{0.46\textwidth}
        \centering
        \includegraphics[width=1\textwidth]{/Users/kavinkannan/Documents/Images/Eigen.png}
        \caption{The Conversion with Meaning\textsuperscript{4}}
        \label{fig:image1}
    \end{minipage} \hfill
\end{figure}
\vspace{1em}
\begingroup
This incredibly arduous process nets us an understanding of the physical world that we 
otherwise wouldn't be able to obtain. Although, definitely not applicable to everybody's 
daily life, the Hamiltonian operator's eigenavalues and eigenvectors are extremely important 
in the context of quantum mechanics, which fundamentally controls our understanding of 
planetary motion, molecular behavior, and chaotic systems. Without the Eigenvalues and 
Eigenvectors, the Hamiltonian couldn't net us the values that control the known universe.\\
\rule{\textwidth}{0.4pt}

\newpage
\thispagestyle{empty}
\begingroup %High Dimensional Matrix
Linear Algebra builds off of itself in increasing magnitudes of difficulty. The most 
elementary concept of linear algebra are matrices, which allow for easier storage and organization
of massive quantities of information. There is another type of matrices, namely "higher-
dimensional matrices" that are used in physics extensively. These matrices allow for the 
understanding and computation of tensors. Tensors, most commonly, are used to describe 
the "stress in a material object"\textsuperscript{5}. Stress is in essence, the amount of force per unit area, 
typically causing deformation.
\endgroup
\vspace{1em}
\begin{figure}[h]
    \centering
    \begin{minipage}{0.46\textwidth}
        \centering
        \includegraphics[width=1\textwidth]{/Users/kavinkannan/Documents/Images/Tensor.png}
        \caption{A "higher-dimensional matrix"\textsuperscript{6}}
        \label{fig:image1}
    \end{minipage} \hfill
\end{figure}
\vspace{1em}\\
\begingroup
This already complicated matrix of matrices is further complicated by the calculations that 
have to be done with it in order to net usable information. Although tensors are often 
unused, their applications go far beyond just theoretical physics. From machine learning to 
describing the curvature of spacetime, the uses of tensors are innumerable with their understanding 
rooted in linear algebra.
\endgroup
\begin{figure}[h]
    \centering
    \begin{minipage}{0.45\textwidth}
        \centering
        \includegraphics[width=1\textwidth]{/Users/kavinkannan/Documents/Images/Machine.png}
        \caption{Tensors in Machine Learning\textsuperscript{7}}
        \label{fig:image1}
    \end{minipage} \hfill
    \begin{minipage}{0.45\textwidth}
        \centering
        \includegraphics[width=0.9\textwidth]{/Users/kavinkannan/Documents/Images/Spacetime.jpg}
        \caption{Tensors for Spacetime\textsuperscript{8}}
        \label{fig:image2}
    \end{minipage}
\end{figure}\\
\rule{\textwidth}{0.4pt}

\begingroup
In essence, linear algebra as a subject of learning is critical to the inner and outer working 
of physics. There are countless other applications of linear algebra throughout physics, 
noticed and unnoticed as linear algebra in it of itself is the basis for almost all higher level 
studies of any topic. The applications of the various topics within linear algebra is almost 
limitless, and I personally can't wait to learn about them.
\endgroup

\newpage
\thispagestyle{empty}
\begingroup %Citations
    \centering
    \LARGE Works Cited\\
\endgroup
\vspace{2em}

\begingroup
1. Kirchhoff's laws. \textit{Isaac Physics}. Available at:\\ 
\href{https://isaacphysics.org/concepts/cp_kirchhoffs_laws}{Issac Physics: Kirchoff's Laws}. Accessed: 28 November 2024.
\vspace{1em}\\
\indent 2. Andrew Craig. Kirchhoff's Laws. \textit{SlidePlayer}. Available at:\\ 
\href{https://slideplayer.com/slide/13365572/}{SlidePlayer: Kirchoff's Laws}. Accessed: 28 November 2024.
\vspace{1em}\\
\indent 3. Hamiltonian. \textit{QuEra}. Available at:\\
\href{https://www.quera.com/glossary/hamiltonian}{QuEra: Hamiltonian}. Accessed: 28 November 2024.
\vspace{1em}\\
\indent 4. Professor Dave. Unpacking the Schrödinger Equation. \textit{YouTube}. Available at:\\ 
\href{https://www.youtube.com/watch?app=desktop&v=l4fFG2utivw}{YouTube: Unpacking the Schrödinger Equation - Professor Dave Explains}. Accessed: 28 November 2024.
\vspace{1em}\\
\indent 5. Joseph C. Kolecki. An Introduction To Tensors for Students of Physics and Engineering. \textit{NASA}. Available at:\\
\href{https://www.grc.nasa.gov/www/k-12/Numbers/Math/documents/Tensors_TM2002211716.pdf}{NASA: An Introduction To Tensors for Students of Physics and Engineering - Joseph C. Kolecki}. Accessed: 28 November 2024.
\vspace{1em}\\
\indent 6. Hunter Phillips. A Simple Introduction to Tensors. \textit{Medium}. Available at:\\
\href{https://medium.com/@hunter-j-phillips/a-simple-introduction-to-tensors-c4a8321efffc}{Medium: A Simple Introduction to Tensors - Hunter Phillips}. Accessed: 29 November 2024.
\vspace{1em}\\
\indent 7. amoeba. Why the sudden fascination with tensors? \textit{StackExchange}. Available at:\\
\href{https://stats.stackexchange.com/questions/198061/why-the-sudden-fascination-with-tensors}{Why the sudden fascination with tensors?}. Accessed: 29 November 2024.
\vspace{1em}\\
\indent 8. Ville Hirvonen. The Metric Tensor: A Complete Guide With Examples. \textit{Profound Physics}. Available at:\\
\href{https://profoundphysics.com/metric-tensor-a-complete-guide-with-examples/}{Profound Physics: The Metric Tensor: A Complete Guide With Examples}. Accessed: 29 November 2024.
\endgroup
\end{document}